% ============================================================================
% CHAPTER 1H: BIOMEDICAL SYSTEMS
% ============================================================================

\chapter{Biomedical Systems}
\label{ch:biomedical}

Biomedical control systems interface with the human body, requiring special consideration for safety, variability between patients, and the inherent complexity of biological processes.

% ============================================================================
\section{Pharmacokinetics: Drug Distribution}
\label{sec:pharmacokinetics}
% ============================================================================

Pharmacokinetics describes how drugs move through the body.

\subsection{One-Compartment Model}

\input{figures/fig_ch14_one_compartment.tex}

Mass balance:
\begin{equation}
V\frac{\dd C}{\dd t} = u(t) - k_e V C
\end{equation}

where:
\begin{itemize}
    \item $V$ = volume of distribution (L)
    \item $C$ = drug concentration (mg/L)
    \item $u(t)$ = infusion rate (mg/min)
    \item $k_e$ = elimination rate constant (1/min)
\end{itemize}

Transfer function:
\begin{equation}
\frac{C(s)}{U(s)} = \frac{1/V}{s + k_e}
\end{equation}

Half-life: $t_{1/2} = \ln(2)/k_e$

\begin{examplebox}[title=Example 1.15: IV Drug Infusion]
A drug has $V = 50$ L and $k_e = 0.1$ min$^{-1}$. Find:
\begin{enumerate}[label=(\alph*)]
    \item The half-life
    \item The infusion rate needed for steady-state concentration of 10 mg/L
    \item Time to reach 90\% of steady-state
\end{enumerate}

\textbf{Solution:}

(a) Half-life: $t_{1/2} = \ln(2)/0.1 = 6.93$ min

(b) At steady-state: $u_{ss} = k_e V C_{ss} = 0.1 \times 50 \times 10 = 50$ mg/min

(c) Time constant: $\tau = 1/k_e = 10$ min

Time to 90\%: $t = -\tau\ln(0.1) = 10 \times 2.3 = 23$ min
\end{examplebox}

\subsection{Two-Compartment Model}

\input{figures/fig_ch14_two_compartment.tex}

Mass balances:
\begin{align}
V_1\frac{\dd C_1}{\dd t} &= u(t) - (k_{10} + k_{12})V_1 C_1 + k_{21}V_2 C_2 \\
V_2\frac{\dd C_2}{\dd t} &= k_{12}V_1 C_1 - k_{21}V_2 C_2
\end{align}

This yields a second-order system with two exponential phases:
\begin{itemize}
    \item \textbf{Distribution phase} ($\alpha$): rapid initial decline
    \item \textbf{Elimination phase} ($\beta$): slower terminal decline
\end{itemize}

% ============================================================================
\section{Glucose-Insulin Dynamics}
\label{sec:glucose}
% ============================================================================

\input{figures/fig_ch14_glucose_insulin.tex}

\subsection{Minimal Model (Bergman)}

\begin{align}
\frac{\dd G}{\dd t} &= -p_1 G - X(G + G_b) + D(t) \\
\frac{\dd X}{\dd t} &= -p_2 X + p_3(I - I_b) \\
\frac{\dd I}{\dd t} &= -n(I - I_b) + \gamma(G - h)^+ + u(t)
\end{align}

where:
\begin{itemize}
    \item $G$ = glucose concentration (mg/dL)
    \item $I$ = insulin concentration ($\mu$U/mL)
    \item $X$ = insulin action (remote compartment)
    \item $D(t)$ = meal disturbance
    \item $u(t)$ = insulin infusion (for artificial pancreas)
\end{itemize}

\subsection{Linearized Model for Control Design}

Near operating point $(G_0, I_0)$:
\begin{equation}
\begin{bmatrix} \Delta\dot{G} \\ \Delta\dot{I} \end{bmatrix} =
\begin{bmatrix} -p_1 - X_0 & -K_G \\ K_I & -n \end{bmatrix}
\begin{bmatrix} \Delta G \\ \Delta I \end{bmatrix} +
\begin{bmatrix} 1 \\ 0 \end{bmatrix} \Delta D +
\begin{bmatrix} 0 \\ 1 \end{bmatrix} \Delta u
\end{equation}

% ============================================================================
\section{Cardiovascular System}
\label{sec:cardiovascular}
% ============================================================================

\subsection{Windkessel Model}

The arterial system can be modeled as an electrical analog:

\input{figures/fig_ch14_windkessel.tex}

Governing equation:
\begin{equation}
C\frac{\dd P}{\dd t} + \frac{P}{R} = Q(t)
\end{equation}

where:
\begin{itemize}
    \item $P$ = arterial pressure (mmHg)
    \item $Q$ = cardiac output (mL/s)
    \item $C$ = arterial compliance (mL/mmHg)
    \item $R$ = systemic vascular resistance (mmHg$\cdot$s/mL)
\end{itemize}

\subsection{Blood Pressure Control}

The baroreflex regulates blood pressure:
\begin{equation}
\frac{\dd P}{\dd t} = -\frac{1}{\tau_P}(P - P_{\text{set}}) + K_u u(t)
\end{equation}

where $u(t)$ represents vasoactive drug infusion.

% ============================================================================
\section{Respiratory System}
\label{sec:respiratory}
% ============================================================================

\subsection{Mechanical Ventilation Model}

\input{figures/fig_ch14_lung_model.tex}

Equation of motion:
\begin{equation}
P_{aw} = R_{aw}\dot{V} + \frac{V}{C_L} + P_{\text{PEEP}}
\end{equation}

where:
\begin{itemize}
    \item $P_{aw}$ = airway pressure (cmH$_2$O)
    \item $R_{aw}$ = airway resistance (cmH$_2$O$\cdot$s/L)
    \item $C_L$ = lung compliance (L/cmH$_2$O)
    \item $V$ = lung volume (L)
    \item $P_{\text{PEEP}}$ = positive end-expiratory pressure
\end{itemize}

Transfer function:
\begin{equation}
\frac{V(s)}{P_{aw}(s)} = \frac{C_L}{R_{aw}C_L s + 1}
\end{equation}

% ============================================================================
\section{Anesthesia Delivery}
\label{sec:anesthesia}
% ============================================================================

\subsection{Propofol Pharmacokinetics}

Three-compartment mammillary model:
\begin{align}
V_1\dot{C}_1 &= u - k_{10}V_1 C_1 - k_{12}V_1 C_1 + k_{21}V_2 C_2 - k_{13}V_1 C_1 + k_{31}V_3 C_3 \\
V_2\dot{C}_2 &= k_{12}V_1 C_1 - k_{21}V_2 C_2 \\
V_3\dot{C}_3 &= k_{13}V_1 C_1 - k_{31}V_3 C_3
\end{align}

\subsection{Pharmacodynamic Effect}

The effect-site concentration $C_e$ relates to the measured output (e.g., BIS index):
\begin{equation}
\frac{\dd C_e}{\dd t} = k_{e0}(C_1 - C_e)
\end{equation}

Hill equation for effect:
\begin{equation}
\text{Effect} = E_0 - E_{\max}\frac{C_e^\gamma}{C_e^\gamma + EC_{50}^\gamma}
\end{equation}

% ============================================================================
\section{Medical Device Examples}
\label{sec:medical_devices}
% ============================================================================

\subsection{Infusion Pump}

Transfer function (actuator dynamics):
\begin{equation}
\frac{Q(s)}{V_{\text{cmd}}(s)} = \frac{K_p}{\tau_p s + 1}
\end{equation}

where $Q$ is flow rate and $V_{\text{cmd}}$ is the command voltage.

\subsection{Pacemaker Rate Response}

Activity-based pacing rate:
\begin{equation}
\tau\frac{\dd R}{\dd t} + R = R_{\text{base}} + K \cdot A(t)
\end{equation}

where $A(t)$ is the activity sensor signal and $R$ is the pacing rate.

\subsection{Cochlear Implant Signal Processing}

Filter bank model:
\begin{equation}
H_k(s) = \frac{\omega_k^2}{s^2 + 2\zeta\omega_k s + \omega_k^2}
\end{equation}

for frequency band $k$ centered at $\omega_k$.

% ============================================================================
\section{Neural Systems}
\label{sec:neural}
% ============================================================================

\subsection{Hodgkin-Huxley Neuron Model}

\begin{equation}
C_m\frac{\dd V}{\dd t} = I_{\text{ext}} - g_{Na}m^3h(V - E_{Na}) - g_K n^4(V - E_K) - g_L(V - E_L)
\end{equation}

with gating variables $m$, $h$, $n$ following first-order kinetics.

\subsection{Simplified Integrate-and-Fire Model}

\begin{equation}
\tau_m\frac{\dd V}{\dd t} = -(V - V_{\text{rest}}) + R_m I_{\text{ext}}
\end{equation}

When $V$ reaches threshold $V_{th}$, the neuron fires and resets.

% ============================================================================
